\subsection{(4) Dataset Statistics}

% 多正样本按相同 description 聚合；表中的 Number 表示图文记录数，而不是文本查询数。
% 当同一文本对应多个图像时，每个图像均计为一条图文记录，因此各行之和为全集记录总数 4,028。
% 难负样本文件未提供 same theme/different object、same object/different attribute、
% same composition/different location 三种类型的独立标签，因此这里仅报告每个查询的难负样本数量分布；
% 三种语义类型可在补充标注后进一步细分。
\begin{table}[htbp]
    \centering
    \begin{minipage}[t]{0.47\textwidth}
        \centering
        \textbf{Multi-Positive Statistics.}
        \label{tab:multi-positive-statistics}
        \begin{tabular}{|l|c|}
            \hline
            Positive images per text & Number \\
            \hline
            1 & 2,904 \\
            2 & 788 \\
            3--5 & 197 \\
            $>5$ & 139 \\
            \hline
            Total & 4,028 \\
            \hline
        \end{tabular}
    \end{minipage}
    \hfill
    \begin{minipage}[t]{0.47\textwidth}
        \centering
        \textbf{Hard Negative Statistics.}
        \label{tab:hard-negative-statistics}
        \begin{tabular}{|l|c|}
            \hline
            Hard negatives per query & Number of queries \\
            \hline
            0 & 3,440 \\
            1--3 & 588 \\
            $>3$ & 0 \\
            \hline
            Total & 4,028 \\
            \hline
        \end{tabular}
    \end{minipage}
\end{table}

% 第4步工作：五个类别的样本统计及其在多正样本和难负样本中的分布（正文文字，后续翻译为英文）。
% 云冈石窟全集共包含 4,028 条图文记录，按照语义内容划分为人物造像、建筑空间、佛教故事、出土文物和装饰纹样五个类别。
% 其中，人物造像有 2,140 条，建筑空间有 1,437 条，佛教故事有 213 条，出土文物有 174 条，装饰纹样有 64 条。
% 多正样本子集共包含 1,102 条记录，主要集中于人物造像和建筑空间两类，分别为 777 条和 280 条；佛教故事、装饰纹样分别为 21 条和 24 条，出土文物类别未出现多正样本。
% 难负样本子集共包含 588 条记录，其中人物造像 340 条、建筑空间 218 条、佛教故事 19 条、装饰纹样 11 条，出土文物类别未出现难负样本。
% 总体而言，多正样本和难负样本主要分布在人物造像与建筑空间类别中，这与全集中两类样本占比较高以及其视觉内容具有较强细粒度相似性有关。
